\section{Figures and Tables}\label{sec:figures}

Figures and tables are where \LaTeX{} and Word differ most. \LaTeX{} moves a float to
the top of a page, the bottom of a page or a page of its own; Word has nothing comparable. \texdocx{} therefore writes every \env{figure} and \env{table} where it
appears in the source, accepts and ignores the placement specifier (\texttt{[htbp]} and
friends), and marks every paragraph of the float before its last caption as
\emph{keep with next}, so an image never ends up on one page with its caption on the
next. What you get is a document a co-author can edit without fighting anchored frames.

\subsection{Images and how they are sized}\label{sec:figures-sizing}

\cmd{includegraphics} follows the \textsf{graphicx} rules, so an image is as large in Word
as it would be in the PDF. The starting point is the image's natural size, computed from
its pixel dimensions and the resolution stored in the file: the \texttt{pHYs} chunk of a
PNG or the density in a JPEG's JFIF header, and 72\,dpi when the file records none,
which is what pdf\TeX{} assumes too. An SVG's natural size comes from its \texttt{width}
and \texttt{height} attributes or, failing those, its \texttt{viewBox}. The options then
apply in the usual way:
\begin{itemize}
  \item \texttt{scale} multiplies the natural size;
  \item \texttt{width} or \texttt{height} alone sets that dimension and scales the other
        to preserve the aspect ratio (\texttt{totalheight} is treated like \texttt{height});
  \item \texttt{width} and \texttt{height} together stretch the image to exactly that box,
        unless \texttt{keepaspectratio} is given, in which case it is scaled to the
        largest size that fits inside the box.
\end{itemize}
Lengths may use any unit and any length register: \cmd{textwidth} and \cmd{linewidth} are
derived from the page set up with \textsf{geometry}, and inside a \env{subfigure} or
\env{minipage} \cmd{linewidth} is the width of that box. After sizing, one rule is
specific to \texdocx: an image wider than the text block is scaled down, keeping its
aspect ratio, to the text width. \LaTeX{} would let it run into the margin with an
\emph{overfull box} warning; Word would let it run off the page. Give the size to
\cmd{includegraphics} itself: \cmd{scalebox} and \cmd{resizebox} around an image are
currently ignored, so the image keeps its own size, and no warning is printed.

Figure~\ref{fig:figures-bars} was included with \verb|width=0.7\linewidth|. Its file is
600 pixels wide and records 150\,dpi, so its natural width is 4\,inches (about
10.2\,cm); the explicit width scales it to 70\,\% of the 16\,cm line, 11.2\,cm, exactly
as in the PDF.

\begin{figure}[htbp]
  \centering
  \includegraphics[width=0.7\linewidth]{bars}
  \caption{A bar chart included at 70\,\% of the line width.}\label{fig:figures-bars}
\end{figure}

Images are looked up in the project folder (by default the folder of the main file) and
in every directory named by
\cmd{graphicspath}; a name without an extension is tried as \texttt{.png}, \texttt{.jpg},
\texttt{.jpeg}, \texttt{.gif}, \texttt{.svg}, \texttt{.bmp}, \texttt{.pdf} and
\texttt{.eps}, in that order. Like every other file a document reads, images must lie
inside the project folder (Section~\ref{sec:security}). PNG, JPEG, GIF, BMP and SVG are
embedded as they are; SVG images carry a blank PNG fallback, because \texdocx{} bundles no
rasteriser, so they display in Word 2016 and later. Word cannot embed PDF or EPS
graphics, so for those \texdocx{} looks for a PNG, SVG or JPEG \emph{twin} with the same
base name and uses it instead; it never runs an external converter.
Table~\ref{tab:figures-formats} summarises the rules. Rotation (\texttt{angle}) and
cropping (\texttt{trim}, \texttt{viewport}) are not applied yet and produce warning
\texttt{W014}.

\begin{table}[htbp]
  \centering
  \caption{How each image format is sized and embedded.}\label{tab:figures-formats}
  \begin{tabular}{@{}lll@{}}
    \toprule
    & \multicolumn{2}{c}{Handling} \\
    \cmidrule(l){2-3}
    Format & Natural size from & Embedded as \\
    \midrule
    PNG & pixels and \texttt{pHYs} resolution & the PNG itself \\
    JPEG & pixels and JFIF density & the JPEG itself \\
    GIF, BMP & pixels at 96\,dpi (BMP: stored resolution) & the file itself \\
    SVG & \texttt{width}/\texttt{height} or \texttt{viewBox} & SVG with blank PNG fallback \\
    \multirow{2}{*}{PDF, EPS} & twin found: the twin's size & the twin \\
    & no twin & placeholder and \texttt{W014} \\
    \bottomrule
  \end{tabular}
\end{table}

\subsection{Subfigures}\label{sec:figures-subfigures}

The \textsf{subcaption} package's \env{subfigure} environment works as in \LaTeX: each
panel gets its own caption, lettered (a), (b), \dots, and a \cmd{label} placed after a
subcaption refers to the panel. Figure~\ref{fig:figures-pair} contains two panels: the
vector diagram in Figure~\ref{fig:figures-diagram} and the photograph in
Figure~\ref{fig:figures-photo}. As in \LaTeX, a reference to a panel combines the figure
number with the panel letter. Two differences remain. Panels are stacked vertically
rather than placed side by side, so leave out the \cmd{hfill} that usually separates
them; \texdocx{} would only approximate it with a tab and warn. And a reference to a panel
is written as plain text such as ``2a'' rather than as a field, so it is not renumbered
when fields are updated in Word.

\begin{figure}[htbp]
  \centering
  \begin{subfigure}{0.45\linewidth}
    \centering
    \includegraphics[width=\linewidth]{diagram.svg}
    \caption{A vector diagram (SVG).}\label{fig:figures-diagram}
  \end{subfigure}
  \begin{subfigure}{0.45\linewidth}
    \centering
    \includegraphics[width=0.8\linewidth]{photo}
    \caption{A photograph (JPEG).}\label{fig:figures-photo}
  \end{subfigure}
  \caption{Two panels in one figure.}\label{fig:figures-pair}
\end{figure}

\subsection{Captions, references and lists}\label{sec:figures-captions}

A caption's number is not typed into the document as plain text. ``Figure~1'' is the word
\emph{Figure}, a non-breaking space and the Word field \verb|SEQ Figure \* ARABIC|, with
the number \LaTeX{} would print stored as the field's cached result; tables use
\verb|SEQ Table|. A \cmd{label} after the caption becomes a bookmark around that field, and
\cmd{ref} becomes a \texttt{REF} field pointing to it, so if a figure is added or deleted
in Word, updating fields (select all, then \textsf{F9}) renumbers the captions and every
reference to them, except the panel references described above. A table caption may come above the table, as in
Table~\ref{tab:figures-formats}; it is then kept with the table that follows it.

For the same reason, \cmd{listoffigures} and \cmd{listoftables} do not produce a static
list. Each inserts a heading and a Word table-of-contents field built from the caption
fields, \verb|TOC \h \z \c "Figure"| or \verb|TOC \h \z \c "Table"|, which Word fills in
the next time fields are updated. This guide does not include one, but a document that
does gets clickable entries with page numbers computed by Word's own layout.

\begin{tip}
Give every float a \cmd{label} placed \emph{after} its \cmd{caption}. A label placed
before the caption picks up the number of the current section or subsection
instead, in \LaTeX{} and in \texdocx{} alike.
\end{tip}
